\chapter{Electroweak Unification}
\label{ch:electroweak-unification}

\section{Overview and strategy}
The goal of electroweak unification is to build a \emph{single} local gauge theory that:
\begin{itemize}
  \item reproduces the observed \textbf{charged-current} weak interactions (mediated by $W^\pm$),
  \item reproduces \textbf{electromagnetism} (mediated by the photon $A_\mu$),
  \item reproduces the observed \textbf{neutral-current} weak interactions (mediated by $Z_\mu$),
  \item and explains why weak interactions are \textbf{short range} (massive vector bosons)
        while electromagnetism is \textbf{long range} (massless photon),
  \item all while keeping the \textbf{local gauge symmetry} exact in the Lagrangian.
\end{itemize}
We will proceed in three big steps:
\begin{enumerate}
  \item Identify the observed currents ($J^\mu_{\rm cc}$, $J^\mu_{\rm em}$, $J^\mu_{\rm nc}$) and rewrite them in a form that suggests a gauge structure.
  \item Build a Yang--Mills theory invariant under $SU_L(2)\times U_Y(1)$ and show how $A_\mu$, $W^\pm_\mu$, $Z_\mu$ emerge after mixing.
  \item Implement \textbf{spontaneous symmetry breaking} via the Higgs mechanism:
        \[
        SU_L(2)\times U_Y(1)\longrightarrow U_{\rm em}(1),
        \]
        generating masses for $W^\pm$ and $Z$ while keeping the photon massless.
\end{enumerate}

\section{Weinberg--Salam model for a single family}
\subsection{Fermion content and observed currents}
We start with one family of leptons and quarks organized into left-handed doublets:
\begin{equation}
L=\begin{pmatrix}\nu_e\\ e\end{pmatrix},
\qquad
Q=\begin{pmatrix}u\\ d\end{pmatrix}.
\end{equation}
The experimentally observed weak charged current (V--A structure) reads
\begin{equation}
\boxed{
J^\mu_{\rm cc}=\bar\nu_e\gamma^\mu P_L e+\bar u\gamma^\mu P_L d
=\sum_{F=L,Q}\bar F\,\tau^+\,\gamma^\mu P_L F,
}
\label{eq:Jcc}
\end{equation}
where $P_L=\frac{1}{2}(1-\gamma_5)$ and $\tau^+$ is the isospin raising operator acting in the doublet space.

The electromagnetic current is vector-like (same coupling to left and right components):
\begin{equation}
\boxed{
J^\mu_{\rm em}=\sum_f Q_f\,\bar f\gamma^\mu f,
}
\label{eq:Je}
\end{equation}
where the sum runs over the fermion fields $f=\nu_e,e,u,d$ (and later also right-handed singlets), and $Q_f$ is the electric charge (in units of $e$).

Finally, the weak neutral current can be written in the general form
\begin{equation}
J^\mu_{\rm nc}
=\sum_f \bar f\gamma^\mu\frac{c_V^f-c_A^f\gamma_5}{2}f
\equiv
\sum_f c_L^f\,\bar f\gamma^\mu P_L f+\sum_f c_R^f\,\bar f\gamma^\mu P_R f.
\end{equation}
The coefficients satisfy
\begin{equation}
c_V^f=c_A^f-2Q_f\,x,
\qquad x\simeq 0.23,
\qquad
c_A^\nu=c_A^u=\frac{1}{2},\quad c_A^e=c_A^d=-\frac{1}{2},
\end{equation}
and therefore
\begin{equation}
\boxed{
c_L^f=c_A^f-Q_f x,\qquad c_R^f=-Q_f x.
}
\label{eq:cLcR}
\end{equation}
An immediate consequence is:
\begin{equation}
Q_f=0\ \Longrightarrow\ c_R^f=0
\quad\Rightarrow\quad
\text{no right-handed neutral-current coupling.}
\end{equation}
In particular, a \emph{right-handed neutrino} (if introduced) would not interact through these currents.

\subsection{A key identity: rewriting the neutral current}
A crucial observation is that the weak neutral current can be expressed as
\begin{equation}
\boxed{
J^\mu_{\rm nc}=\sum_{F=L,Q}\bar F\,\frac{\tau_3}{2}\gamma^\mu P_LF
-x\,J^\mu_{\rm em}.
}
\label{eq:Jnc-T3-Jem}
\end{equation}
This form is extremely suggestive: it separates a purely left-handed isospin piece (built from $\tau_3/2$) and a piece proportional to the electromagnetic current. This is precisely the structure that will appear after mixing two gauge fields into $A_\mu$ and $Z_\mu$.

\subsection{Hypercharge and the electromagnetic current}
Let us focus on the \emph{left-handed} part of the electromagnetic current. One can show
\begin{equation}
\sum_f Q_f\,\bar f\gamma^\mu P_L f
=
\sum_F\left(
\bar F\,\frac{\tau_3}{2}\gamma^\mu P_LF
+\bar F\,\frac{Y_F}{2}\gamma^\mu P_LF
\right),
\end{equation}
where $F$ runs over the left-handed doublets $L$ and $Q$, and $Y_F$ is a new quantum number called \textbf{weak hypercharge}.
For the two left-handed doublets in one family,
\begin{equation}
\boxed{
\frac{Y_L}{2}=-\frac{1}{2},
\qquad
\frac{Y_Q}{2}=\frac{1}{6}.
}
\label{eq:YL-YQ}
\end{equation}
The key point is that the hypercharge term is proportional to the identity in isospin space, so it commutes with $\tau_3$ and with the charged current structure.

Including right-handed singlets, the full electromagnetic current becomes
\begin{equation}
J^\mu_{\rm em}
=
\sum_F \bar F\,\frac{\tau_3}{2}\gamma^\mu P_LF
+
\sum_F \bar F\,\frac{Y_F}{2}\gamma^\mu P_LF
+
\sum_f Q_f\,\bar f\gamma^\mu P_R f.
\end{equation}
From this, the neutral current can be reorganized into a form that separates:
\begin{itemize}
  \item an isospin term,
  \item and a bracket containing the hypercharge (left-handed) and the right-handed charge terms.
\end{itemize}
Explicitly,
\begin{equation}
J^\mu_{\rm nc}
=
(1-x)\sum_F \bar F\,\frac{\tau_3}{2}\gamma^\mu P_LF
-
x\left(
\sum_F \bar F\,\frac{Y_F}{2}\gamma^\mu P_LF
+
\sum_f Q_f\,\bar f\gamma^\mu P_R f
\right).
\end{equation}

\section{Quantum numbers: $T_3$, $Q$ and weak hypercharge}
\subsection{Definition and table}
We choose $SU_L(2)$ weak isospin with generator
\[
T_3=\frac{\tau_3}{2}.
\]
Then the defining relation among electric charge, isospin and hypercharge is
\begin{equation}
\boxed{
Q = T_3 + \frac{Y}{2},
\qquad\text{equivalently}\qquad
\frac{Y}{2}=Q-T_3.
}
\label{eq:Q-T3-Y}
\end{equation}
For one family, the quantum numbers are:

\begin{center}
\begin{tabular}{lcccc}
\hline
Field & $SU_L(2)$ & $T_3$ & $Q_f$ & $\dfrac{Y}{2}\equiv Q_f-T_3$ \\
\hline
$\nu_{eL}$ & $\frac{1}{2}$ & $+\frac{1}{2}$ & $0$ & $-\frac{1}{2}$ \\
$e_L$      & $\frac{1}{2}$ & $-\frac{1}{2}$ & $-1$ & $-\frac{1}{2}$ \\
\hline
$u_L$      & $\frac{1}{2}$ & $+\frac{1}{2}$ & $\frac{2}{3}$ & $\frac{1}{6}$ \\
$d_L$      & $\frac{1}{2}$ & $-\frac{1}{2}$ & $-\frac{1}{3}$ & $\frac{1}{6}$ \\
\hline
$u_R$      & $0$ & $0$ & $\frac{2}{3}$ & $\frac{2}{3}$ \\
$d_R$      & $0$ & $0$ & $-\frac{1}{3}$ & $-\frac{1}{3}$ \\
$e_R$      & $0$ & $0$ & $-1$ & $-1$ \\
\hline
\end{tabular}
\end{center}

\noindent
These assignments will be exactly the ones required for a consistent gauge theory under $SU_L(2)\times U_Y(1)$.

\section{Building the electroweak gauge theory: $SU_L(2)\times U_Y(1)$}
\subsection{Yang--Mills Lagrangian and covariant derivatives}
We now construct a local gauge theory invariant under
\[
SU_L(2)\times U_Y(1).
\]
Introduce gauge fields:
\begin{itemize}
  \item $W_\mu=W_\mu^a T^a$ for $SU_L(2)$, with $T^a=\tau^a/2$,
  \item $B_\mu$ for $U_Y(1)$.
\end{itemize}
The Lagrangian for fermion kinetic terms plus gauge-field kinetic terms is
\begin{equation}
\boxed{
\mathcal L
=
\sum_{F=L_L,Q_L}\bar F\,i\slashed{D}_F F
+\sum_{f=u_R,d_R,e_R}\bar f\,i\slashed{D}_f f
-\frac{1}{2}\mathrm{tr}(W_{\mu\nu}W^{\mu\nu})
-\frac{1}{4}B_{\mu\nu}B^{\mu\nu}.
}
\label{eq:LEW-gauge}
\end{equation}
The covariant derivatives are:
\begin{equation}
\boxed{
D_\mu^F = \partial_\mu + ig W_\mu + ig'\frac{Y_F}{2}B_\mu,
\qquad
D_\mu^f = \partial_\mu + ig'\frac{Y_f}{2}B_\mu,
}
\label{eq:cov-der}
\end{equation}
where $g$ and $g'$ are the electroweak couplings.

The field strengths are:
\begin{align}
W_{\mu\nu}&=T^a W^a_{\mu\nu},\qquad
W^a_{\mu\nu}=\partial_\mu W^a_\nu-\partial_\nu W^a_\mu-g\,\epsilon^{abc}W^b_\mu W^c_\nu,\\
B_{\mu\nu}&=\partial_\mu B_\nu-\partial_\nu B_\mu,
\end{align}
with $[T^a,T^b]=i\epsilon^{abc}T^c$.

\subsection{Gauge transformations}
Under $SU_L(2)\times U_Y(1)$:
\begin{align}
F(x)&\to e^{ig'\frac{Y_F}{2}\theta(x)}\,g_L(x)\,F(x),\\
f(x)&\to e^{ig'\frac{Y_f}{2}\theta(x)}\,f(x),
\end{align}
and the gauge fields transform as
\begin{align}
W_\mu(x)&\to g_L(x)W_\mu(x)g_L^\dagger(x)-\frac{i}{g}\,g_L(x)\partial_\mu g_L^\dagger(x),\\
B_\mu(x)&\to B_\mu(x)-\partial_\mu\theta(x),
\end{align}
which implies
\[
W_{\mu\nu}\to g_L W_{\mu\nu} g_L^\dagger,
\qquad
B_{\mu\nu}\to B_{\mu\nu}.
\]
Thus the Lagrangian \eqref{eq:LEW-gauge} is gauge invariant.

\section{Physical gauge fields and currents}
\subsection{Charged fields $W^\pm_\mu$ and neutral mixing}
Define the charged combinations
\begin{equation}
\boxed{
W^\pm_\mu=\frac{1}{\sqrt{2}}\left(W^1_\mu\mp iW^2_\mu\right).
}
\label{eq:Wpm}
\end{equation}
The neutral fields $B_\mu$ and $W^3_\mu$ mix into the physical photon $A_\mu$ and $Z_\mu$ through the Weinberg angle $\theta_W$:
\begin{align}
A_\mu &= \cos\theta_W\,B_\mu+\sin\theta_W\,W^3_\mu,\\
Z_\mu &= -\sin\theta_W\,B_\mu+\cos\theta_W\,W^3_\mu,
\end{align}
or inverted,
\begin{align}
B_\mu &= \cos\theta_W\,A_\mu-\sin\theta_W\,Z_\mu,\\
W^3_\mu &= \sin\theta_W\,A_\mu+\cos\theta_W\,Z_\mu,
\end{align}
with $s_W\equiv\sin\theta_W$, $c_W\equiv\cos\theta_W$.

\subsection{Charged current interaction}
From the covariant derivative, the charged-current interaction is
\begin{align}
\mathcal L_{\rm cc}
&= -\sum_{F=L_L,Q_L} g\,\bar F\gamma^\mu(T^1W^1_\mu+T^2W^2_\mu)F\\
&= -\sum_{F=L_L,Q_L}\frac{g}{\sqrt{2}}\bar F\gamma^\mu(\tau^+W^+_\mu+\tau^-W^-_\mu)F\\
&=
\boxed{
-\frac{g}{\sqrt{2}}\left(J^\mu_{\rm cc}W^+_\mu+J_{\rm cc}^{\dagger\mu}W^-_\mu\right).
}
\label{eq:Lcc}
\end{align}

\subsection{Electromagnetism and the fundamental relation $gs_W=g'c_W=e$}
The electromagnetic interaction comes from the $A_\mu$ piece inside $W^3_\mu$ and $B_\mu$.
Requiring that \textbf{left and right} components couple in the \emph{same way} to $A_\mu$ enforces:
\begin{equation}
eQ_f = g'\frac{Y_f}{2}c_W = gT_3 s_W + g'\frac{Y_F}{2}c_W.
\end{equation}
Using $\frac{Y}{2}=Q-T_3$, one arrives at the central electroweak relation
\begin{equation}
\boxed{
g\,s_W=g'\,c_W=e.
}
\label{eq:gsW-e}
\end{equation}
This ensures that the photon couples universally to electric charge and that QED emerges as the unbroken gauge symmetry after symmetry breaking.

\subsection{Neutral current interaction and $g_Z$}
Inserting the $Z_\mu$ combinations yields the $Z$ interaction in the compact form
\begin{equation}
\mathcal L_{\rm nc} = -g_Z\,Z_\mu\,J^\mu_{\rm nc},
\qquad
\boxed{
g_Z\equiv \frac{e}{c_W s_W}.
}
\label{eq:gZ}
\end{equation}
The neutral current can be written as
\begin{equation}
\boxed{
J^\mu_{\rm nc}
=
\sum_{F=L_L,Q_L}\bar F\gamma^\mu T_3 F
-
s_W^2 \sum_{f=u,d,e}\,Q_f\,\bar f\gamma^\mu f,
}
\label{eq:Jnc-final}
\end{equation}
which matches the previously identified structure \eqref{eq:Jnc-T3-Jem} with $x=s_W^2$.

\section{The Higgs mechanism}
\subsection{The problem: massless gauge bosons}
The gauge-invariant Lagrangian constructed so far contains only \emph{massless} vector bosons. This would imply long-range weak forces, contradicting observation: weak interactions are extremely short range.

Adding explicit mass terms ``by hand'' is \emph{not allowed} because such terms would break the local gauge symmetry $SU_L(2)\times U_Y(1)$.

\subsection{The solution: spontaneous symmetry breaking}
We use spontaneous symmetry breaking: the Lagrangian is fully symmetric, but the vacuum is not. The mechanism requires a scalar multiplet with a non-zero vacuum expectation value (VEV).

We need at least one \emph{neutral} component to acquire a VEV. Since
\[
Q=T_3+\frac{Y}{2},
\]
a neutral component must satisfy $Q=0$. Taking a scalar doublet with $Y/2=+1/2$ achieves this:
\begin{equation}
\boxed{
\phi=
\begin{pmatrix}
\phi^+\\
\phi^0
\end{pmatrix},
\qquad
\frac{Y_\phi}{2}=\frac{1}{2}.
}
\label{eq:Higgs-doublet}
\end{equation}
Its gauge transformation is
\begin{equation}
\phi(x)\ \longrightarrow\ g_L(x)\,e^{ig'\theta(x)/2}\,\phi(x).
\end{equation}

\subsection{Vacuum structure and unitary gauge}
Assume a potential $V(\phi^\dagger\phi)$ with a minimum at
\begin{equation}
\boxed{
\phi^\dagger\phi=\frac{v^2}{2}\in\mathbb R_+.
}
\label{eq:vev}
\end{equation}
A general field configuration near the minimum can be parameterized as
\begin{equation}
\phi(x) = U(x)\,\frac{1}{\sqrt{2}}
\begin{pmatrix}
0\\
v+h(x)
\end{pmatrix},
\qquad U(x)\in SU(2).
\label{eq:phi-U}
\end{equation}
Important points:
\begin{itemize}
  \item The complex doublet $\phi$ has \textbf{4 real degrees of freedom}.
  \item $h(x)$ is a real field describing fluctuations \emph{along} the radial direction.
  \item $U(x)$ contains the would-be Goldstone bosons if the symmetry were global.
  \item Since the symmetry is \emph{local}, $U(x)$ can be removed by a gauge transformation:
        this is the \textbf{unitary gauge}.
\end{itemize}
In unitary gauge,
\begin{equation}
\boxed{
\phi(x)=\frac{1}{\sqrt{2}}
\begin{pmatrix}
0\\
v+h(x)
\end{pmatrix}.
}
\label{eq:unitary-gauge}
\end{equation}

\subsection{Unbroken subgroup and breaking pattern}
The vacuum configuration $h(x)=0$ is invariant under transformations with
\[
g_L(x)=e^{ig'\theta(x)T_3},
\]
which corresponds precisely to the electromagnetic subgroup.
Hence the breaking pattern is:
\begin{equation}
\boxed{
SU_L(2)\times U_Y(1)\ \longrightarrow\ U_{\rm em}(1).
}
\label{eq:SSB-pattern}
\end{equation}

\subsection{Gauge boson masses from the Higgs kinetic term}
The scalar kinetic term is
\begin{equation}
(D_\mu\phi)^\dagger D^\mu\phi,
\qquad
D_\mu=\partial_\mu+igW_\mu+ig'\frac{1}{2}B_\mu.
\end{equation}
In terms of physical fields it can be written (as given) as
\begin{equation}
\begingroup\setlength{\jot}{2pt}
\begin{aligned}
D_\mu
&=
\partial_\mu
+i\frac{g}{\sqrt{2}}(\tau^+W^+_\mu+\tau^-W^-_\mu)
+igT_3 W^3_\mu+ig'\frac{1}{2}B_\mu
\\
&=
\partial_\mu
+i\frac{g}{\sqrt{2}}(\tau^+W^+_\mu+\tau^-W^-_\mu)
+\begin{pmatrix}
ieA_\mu+ig_Z(c_W^2-s_W^2)Z_\mu & 0\\
0 & -ig_Z\frac{1}{2}Z_\mu
\end{pmatrix}.
\end{aligned}
\endgroup
\end{equation}

Now evaluate at the vacuum configuration ($h(x)=0$):
\[
\phi(x)=\frac{1}{\sqrt{2}}
\begin{pmatrix}
0\\ v
\end{pmatrix}.
\]
Then
\begin{equation}
D_\mu\phi=
\begin{pmatrix}
i\frac{gv}{2}W^+_\mu\\[1mm]
-\;i\frac{g_Z v}{2\sqrt{2}}Z_\mu
\end{pmatrix},
\end{equation}
and therefore
\begin{equation}
(D_\mu\phi)^\dagger D^\mu\phi
=
\frac{g^2v^2}{4}\,W^-_\mu W^{+\mu}
+\frac{g_Z^2v^2}{8}\,Z_\mu Z^\mu.
\end{equation}
This is precisely a mass term for $W^\pm$ and $Z$. Reading off the masses,
\begin{equation}
\boxed{
m_W=\frac{gv}{2},
\qquad
m_Z=\frac{g_Z v}{2},
\qquad
g=g_Z c_W\ \Rightarrow\ m_Z>m_W.
}
\label{eq:mW-mZ}
\end{equation}
The photon remains massless because $U_{\rm em}(1)$ is unbroken.

\subsection{Degrees of freedom reshuffling: ``eating'' Goldstones}
Before symmetry breaking:
\begin{itemize}
  \item 3 massless gauge bosons $\Rightarrow 3\times 2=6$ d.o.f.,
  \item 2 complex scalars $\Rightarrow 2\times 2=4$ d.o.f.,
\end{itemize}
total $10$ degrees of freedom.

After symmetry breaking:
\begin{itemize}
  \item 3 massive gauge bosons ($W^\pm$, $Z$) $\Rightarrow 3\times 3=9$ d.o.f.,
  \item 1 real scalar ($h$) $\Rightarrow 1$ d.o.f.,
\end{itemize}
total again $10$.

Thus, the would-be Goldstone bosons are removed by a gauge transformation and become the longitudinal polarizations of the massive vector bosons: this is the \textbf{Higgs mechanism}. The remaining scalar field $h(x)$ describes a physical neutral scalar particle: the \textbf{Higgs boson}.

\subsection{Historical note (as stated)}
Some physicists speculated the field $h(x)$ could be superfluous (no Higgs particle), while others generalized to several scalar multiplets (several Higgs bosons). The Higgs boson was found at the LHC (CERN) in 2013, and so far no other fundamental scalar particle has showed up.

\section{Matching to Fermi theory and numerical predictions}
At low energy $E\ll m_W,m_Z$, amplitudes must reproduce the effective four-fermion Fermi theory. Matching yields:
\begin{equation}
\boxed{
\frac{g^2}{2m_W^2}=2\sqrt{2}\,G,
\qquad
\frac{g_Z^2}{m_Z^2}=4\sqrt{2}\,\rho\,G.
}
\label{eq:matching}
\end{equation}
From these equalities one obtains
\begin{equation}
\boxed{
\rho=\frac{g_Z^2}{m_Z^2}\frac{m_W^2}{g^2}=1,
\qquad
v=\frac{1}{\sqrt{\sqrt{2}G}}\simeq 246~\mathrm{GeV}.
}
\label{eq:rho-v}
\end{equation}
Using $s_W^2\simeq 0.23$ from low-energy experiments, one gets
\begin{equation}
g=\frac{e}{s_W}\simeq 0.63,
\qquad
g'=\frac{e}{c_W}\simeq 0.34,
\qquad
g_Z=\frac{e}{s_Wc_W}\simeq 0.72,
\end{equation}
and the corresponding (tree-level) mass estimates:
\begin{equation}
m_W\simeq 78~\mathrm{GeV},
\qquad
m_Z\simeq 89~\mathrm{GeV}.
\end{equation}
This constituted a prediction for $m_W$ and $m_Z$, which were found at CERN in 1983.

The current experimental values quoted are:
\begin{equation}
m_W\simeq 80.379(12)~\mathrm{GeV},
\qquad
m_Z\simeq 91.1876(21)~\mathrm{GeV}.
\end{equation}

\section{Fermion masses from Yukawa terms}
\subsection{Why fermions are massless without Yukawas}
So far, we have generated masses for the vector bosons while preserving gauge symmetry in the Lagrangian. However, fermions would remain massless unless we add additional gauge-invariant terms.

\subsection{Gauge-invariant Yukawa interactions}
The scalar doublet $\phi$ allows \textbf{Yukawa terms}:
\begin{equation}
\boxed{
\mathcal L_{\rm Yuk}
=
-\lambda_e\,\bar L\,\phi\,e_R
-\lambda_d\,\bar Q\,\phi\,d_R
-\lambda_u\,\bar Q\,\phi^c\,u_R
+\mathrm{H.c.}
}
\label{eq:Yukawa}
\end{equation}
where
\begin{equation}
\boxed{
\phi^c \equiv i\tau_2\,\phi^\ast,
\qquad
\phi^c(x)\ \longrightarrow\ g_L(x)e^{-i\theta(x)/2}\,\phi(x).
}
\label{eq:phic}
\end{equation}
In unitary gauge,
\[
\phi(x)=\frac{1}{\sqrt{2}}\begin{pmatrix}0\\ v+h(x)\end{pmatrix},
\qquad
\phi^c(x)=\frac{1}{\sqrt{2}}\begin{pmatrix}v+h(x)\\ 0\end{pmatrix}.
\]
Substituting into \eqref{eq:Yukawa} gives
\begin{equation}
\mathcal L_{\rm Yuk}
=
-\left(\lambda_e\,\bar e_L e_R+\lambda_d\,\bar d_L d_R+\lambda_u\,\bar u_L u_R\right)\frac{1}{\sqrt{2}}(v+h)
+\mathrm{H.c.}
\end{equation}
This splits into a mass term and a Higgs--fermion interaction:
\begin{align}
\boxed{
m_f = \frac{\lambda_f v}{\sqrt{2}},
\qquad f=e,d,u,
}
\label{eq:mf-lambda}
\end{align}
and
\begin{equation}
\boxed{
\mathcal L_{h\bar f f}
= -\left(m_e\bar e_L e_R+m_d\bar d_L d_R+m_u\bar u_L u_R+\mathrm{H.c.}\right)\frac{h}{v}.
}
\label{eq:hff}
\end{equation}
Thus, \textbf{the Higgs coupling to a fermion is proportional to the fermion mass}.

\section{Higgs interactions with vector bosons and self-interactions}
\subsection{Vector boson couplings}
From the Higgs kinetic term one obtains, as given,
\begin{equation}
(D_\mu\phi)^\dagger D^\mu\phi
=
\frac{1}{2}\partial_\mu h\,\partial^\mu h
+m_W^2 W^-_\mu W^{+\mu}\left(1+\frac{h}{v}\right)^2
+\frac{m_Z^2}{2}Z_\mu Z^\mu\left(1+\frac{h}{v}\right)^2.
\end{equation}
This explicitly shows:
\begin{itemize}
  \item the Higgs kinetic term $\frac{1}{2}(\partial h)^2$,
  \item the $W$ and $Z$ mass terms (from the $1$ inside $(1+h/v)^2$),
  \item and the Higgs couplings to $W$ and $Z$ (from the linear and quadratic terms in $h/v$).
\end{itemize}
Since $v\simeq 246~\mathrm{GeV}$, these couplings are typically small, except for the top quark ($m_t\simeq 173~\mathrm{GeV}$) and for the heavy gauge bosons $W^\pm$ and $Z^0$.

\subsection{Higgs self-interaction from the potential}
The potential $V(\phi^\dagger\phi)$ with a non-zero minimum produces a Higgs mass and self-interactions.
Using the stated analogy with the linear sigma model, one may write:
\begin{equation}
\boxed{
\mathcal L_h
=
-\frac{m_h^2}{2}h^2
-\frac{m_h^2}{2v}h^3
-\frac{m_h^2}{8v^2}h^4.
}
\label{eq:Higgs-self}
\end{equation}
Here $m_h$ was initially a free parameter (hardly constrained by other observables).
Once $m_h$ is measured, the Higgs self-interactions are fixed.

The quoted current Higgs mass value is:
\begin{equation}
m_h=125.10(14)~\mathrm{GeV}.
\end{equation}

\section{More families and flavor structure}
\subsection{Generalizing the fermion content}
In nature there are at least three families. For an arbitrary number of families $N_f$:
\begin{equation}
L_i=\begin{pmatrix}\nu_{l_i}\\ l_i\end{pmatrix},
\qquad
Q_i=\begin{pmatrix}u_i\\ d_i\end{pmatrix},
\qquad i=1,\dots,N_f,
\end{equation}
with examples:
\[
(L_1,Q_1)=\left(\begin{pmatrix}\nu_e\\ e\end{pmatrix},\begin{pmatrix}u\\ d\end{pmatrix}\right),\quad
(L_2,Q_2)=\left(\begin{pmatrix}\nu_\mu\\ \mu\end{pmatrix},\begin{pmatrix}c\\ s\end{pmatrix}\right),\quad
(L_3,Q_3)=\left(\begin{pmatrix}\nu_\tau\\ \tau\end{pmatrix},\begin{pmatrix}t\\ b\end{pmatrix}\right).
\]
The gauge-field part of the Lagrangian is unchanged. The fermion kinetic terms become
\begin{equation}
\mathcal L_{\rm kin}^{\rm fermions}
=
\sum_{i=1}^{N_f}\left(
\sum_{F_i=L_{iL},Q_{iL}}\bar F_i\,i\slashed{D}_F F_i
+\sum_{f_i=u_{iR},d_{iR},e_{iR}}\bar f_i\,i\slashed{D}_f f_i
\right).
\end{equation}

\subsection{General Yukawa matrices and mass diagonalization}
With several families, Yukawa terms become matrices in family space:
\begin{equation}
\boxed{
\mathcal L_{\rm Yuk}
=
\sum_{i,j=1}^{N_f}\left[
-\lambda^{l}_{ij}\,\bar L_{iL}\phi\,l_{jR}
-\lambda^{d}_{ij}\,\bar Q_{iL}\phi\,d_{jR}
-\lambda^{u}_{ij}\,\bar Q_{iL}\phi^c\,u_{jR}
+\mathrm{H.c.}
\right].
}
\label{eq:Yukawa-matrix}
\end{equation}
This leads to \textbf{off-diagonal mass terms}: fields with well-defined $SU_L(2)\times U_Y(1)$ quantum numbers are not mass eigenstates.

To obtain mass eigenstates, one diagonalizes the mass matrices through \textbf{unitary} field redefinitions (to keep canonical kinetic terms).
For down-type quarks, write the mass term schematically as
\begin{equation}
\bar d_L\,\Lambda_d\,d_R
=\sum_{i,j}\bar d_{iL}\,\Lambda^d_{ij}\,d_{jR},
\qquad
\boxed{
\Lambda^d_{ij}\equiv \lambda^d_{ij}\frac{v}{\sqrt{2}}.
}
\end{equation}
Assuming $\det\Lambda_d\neq 0$, one writes $\Lambda_d$ as a unitary times a Hermitian matrix:
\begin{equation}
\boxed{
\Lambda_d = U_d M_d,
\qquad
M_d=S_d D_d^{1/2}S_d^\dagger,
\qquad
\Lambda_d^\dagger\Lambda_d=S_d D_d S_d^\dagger,
}
\label{eq:Lambda-decomp}
\end{equation}
where $S_d$ is unitary and $D_d$ is diagonal positive definite (from diagonalizing the Hermitian matrix $\Lambda_d^\dagger\Lambda_d$).
One can verify the unitarity of $U_d$ as stated:
\[
U_dU_d^\dagger=\mathbf 1.
\]
Then the down-quark mass term becomes diagonal under
\begin{equation}
\boxed{
d_R\to S_d\,d_R,
\qquad
d_L\to U_d S_d\,d_L,
}
\end{equation}
where $d_R$ and $d_L$ are vectors in family space (indices suppressed).

One proceeds analogously for up-type quarks and charged leptons.

\subsection{Why neutral currents remain diagonal, and charged currents do not}
Because the transformations are unitary, they do \emph{not} affect terms that are diagonal in isospin space, namely:
\begin{itemize}
  \item the electromagnetic current,
  \item the neutral current.
\end{itemize}
However, the charged current involves $u_L$ and $d_L$ together. Since the unitary transformations for up and down quarks need not be the same, the charged current acquires a mixing matrix:
\begin{equation}
\boxed{
\bar u_L\gamma^\mu d_L \ \longrightarrow\ \bar u_L\gamma^\mu\,V_{\rm CKM}\,d_L,
\qquad
V_{\rm CKM}\equiv S_u^\dagger U_u^\dagger U_d S_d.
}
\label{eq:CKM}
\end{equation}
This is the origin of the \textbf{CKM matrix} in the electroweak theory.

\subsection{Leptons and the PMNS comment}
For leptons, the charged current transforms as
\begin{equation}
\bar\nu_{lL}\gamma^\mu l_L \ \longrightarrow\ \bar\nu_{lL}\gamma^\mu U_l S_l\,l_L.
\end{equation}
Since there is no neutrino mass term in this setup, one can perform a unitary transformation on the neutrino fields to diagonalize the charged current.
If one includes a right-handed neutrino, then a neutrino mass term becomes possible and the situation becomes analogous to the quark case: a \textbf{PMNS matrix} arises.

\section{The GIM mechanism and suppression of FCNC}
\subsection{The empirical fact}
Flavor-changing neutral currents (FCNC) are very suppressed. For instance,
\begin{equation}
\frac{\Gamma(K_L^0\to \mu^+\mu^-)}{\Gamma(K_L^0\to \text{all})}\simeq 9\times 10^{-9}.
\end{equation}
A naive expectation from weak interactions might suggest a suppression of order
\[
\alpha_W^2,
\qquad
\alpha_W\equiv \frac{g^2}{4\pi}\simeq 0.032,
\qquad
\alpha_W^2\simeq 10^{-3},
\]
which is \emph{far} too weak compared with $10^{-9}$.

\subsection{Origin: unitarity of the CKM matrix (GIM)}
The strong suppression arises from the unitarity of the CKM matrix: this is the \textbf{GIM mechanism} (Glashow--Iliopoulos--Maiani, 1970).

To illustrate, consider two families with Cabibbo mixing:
\begin{equation}
\begin{pmatrix}d'\\ s'\end{pmatrix}
=
\begin{pmatrix}
\cos\theta_c & \sin\theta_c\\
-\sin\theta_c & \cos\theta_c
\end{pmatrix}
\begin{pmatrix}d\\ s\end{pmatrix}.
\end{equation}
Then the charged current is
\begin{equation}
J^\mu
=
\bar\nu_e\gamma^\mu P_L e+\bar\nu_\mu\gamma^\mu P_L \mu
+\bar u\gamma^\mu P_L d'
+\bar c\gamma^\mu P_L s'.
\end{equation}
For $K_L^0\to \mu^+\mu^-$ there are two loop diagrams contributing (as shown in the slides).
If one neglects $m_u,m_c\ll m_W$, the two contributions cancel.

Including quark masses, the cancellation is not exact and produces an \emph{extra} suppression:
\begin{equation}
\boxed{
\mathcal A_{\rm FCNC}\ \sim\ \frac{m_c^2}{m_W^2}\ \sim\ 3.2\times 10^{-4},
}
\end{equation}
and therefore for the decay width,
\begin{equation}
\boxed{
\Gamma_{\rm FCNC}\ \sim\ \left(\frac{m_c^2}{m_W^2}\right)^2
=\frac{m_c^4}{m_W^4}\ \sim\ 10^{-7}.
}
\end{equation}
This explains why FCNC processes can be far more suppressed than a naive $\alpha_W^2$ estimate.

\subsection{Third family contribution and CKM hierarchy}
Including the third family adds a contribution involving the top quark:
\begin{itemize}
  \item since $m_t>m_W$, it could be potentially large,
  \item but the relevant CKM elements are small:
\[
|V_{ts}|\sim \lambda^2,\qquad |V_{td}|\sim \lambda^3,\qquad \lambda\sim 0.22,
\]
so the contribution is suppressed by
\begin{equation}
\boxed{
\sim \lambda^5 \sim 5\times 10^{-4},
}
\end{equation}
a suppression comparable in size to the two-family charm mechanism.
Similar (not identical) suppressions occur in $D^0\to \mu^+\mu^-$, $B^0\to \mu^+\mu^-$ and non-leptonic FCNC decays.

Since FCNC are very suppressed within the electroweak theory, they provide a good place to search for effects of physics beyond the Standard Model.

\section{Electroweak phenomenology}
\subsection{$W^\pm$ leptonic decays and lepton universality}
Consider
\[
W_A^+\to l_1^+\,\nu_{l\,2}.
\]
The relevant interaction is
\begin{equation}
\mathcal L_{\rm int}= -\frac{g}{\sqrt{2}}\,J^\mu_{\rm cc}W^+_\mu,
\qquad
J^\mu_{\rm cc}\simeq \bar\nu_l\gamma^\mu P_L l.
\end{equation}
The matrix element is
\begin{equation}
\mathcal M
= -\frac{g}{\sqrt{2}}\,
\bar u(2)\gamma^\mu P_L v(1)\,\epsilon_\mu(A).
\end{equation}
Neglecting lepton masses and averaging over $W^+$ polarizations, the partial width is
\begin{equation}
\boxed{
\Gamma(W^+\to l^+\nu_l)
=
\frac{\alpha_W\,m_W}{12}.
}
\label{eq:Wleptonic}
\end{equation}
Numerically this matches the quoted values:
\[
\Gamma_{\rm Exp}\simeq 0.22~\mathrm{GeV},
\qquad
\Gamma_{\rm Th}\simeq 0.21~\mathrm{GeV}.
\]
The result is the same for any light lepton family: \textbf{lepton universality}.

\subsection{$W^\pm$ hadronic decays, CKM suppression, and colors}
Using the duality hypothesis:
\[
\Gamma(W^+\to \text{hadrons})=\Gamma(W^+\to \text{quarks}).
\]
The quark-level channels are
\begin{align}
\Gamma(W^+\to \text{quarks})&=
\Gamma(W^+\to u\bar d)+\Gamma(W^+\to u\bar s)+\Gamma(W^+\to u\bar b)\nonumber\\
&\quad+\Gamma(W^+\to c\bar d)+\Gamma(W^+\to c\bar s)+\Gamma(W^+\to c\bar b).
\end{align}
The CKM hierarchy implies suppressions such as
\[
\Gamma(W^+\to u\bar s)\sim |V_{us}|^2\sim \lambda^2,
\]
\[
\Gamma(W^+\to u\bar b)\sim |V_{ub}|^2\sim \lambda^6,
\]
\[
\Gamma(W^+\to c\bar d)\sim |V_{cd}|^2\sim \lambda^2,
\]
\[
\Gamma(W^+\to c\bar b)\sim |V_{cb}|^2\sim \lambda^4.
\]

The notes emphasize that QCD corrections are of order $\alpha_s(m_Z)\sim 0.11$, larger than $\lambda^2\sim 0.048$.
Therefore, at leading order (neglecting quark masses),
\begin{equation}
\Gamma(W^+\to \text{hadrons})
\simeq \Gamma(W^+\to u\bar d)+\Gamma(W^+\to c\bar s).
\end{equation}
Each quark channel has an extra factor $N_c=3$ (number of colors), so:
\begin{equation}
\boxed{
\Gamma(W^+\to \text{hadrons})
\simeq 2\times N_c \times \Gamma(W^+\to l^+\nu_l)
=
\frac{\alpha_W m_W N_c}{6}.
}
\label{eq:Whadronic}
\end{equation}
Numerically:
\[
\Gamma_{\rm Exp}\simeq 1.4~\mathrm{GeV},
\qquad
\Gamma_{\rm Th}\simeq 1.3~\mathrm{GeV}.
\]

\subsection{$Z^0$ decays: couplings and partial widths}
The interaction is
\begin{equation}
\mathcal L_{\rm int}=-g_Z\,J^\mu_{\rm nc}Z_\mu,
\qquad g_Z=\frac{g}{c_W}.
\end{equation}
For each fermion:
\begin{equation}
J^\mu_{\rm nc}\simeq
C_L^f\,\bar f\gamma^\mu P_L f
+
C_R^f\,\bar f\gamma^\mu P_R f,
\end{equation}
with
\begin{equation}
\boxed{
C_L^f=T_3-Q_f s_W^2,
\qquad
C_R^f=-Q_f s_W^2.
}
\label{eq:CLCR}
\end{equation}
Since $m_f\ll m_Z$ for all fermions except the top quark ($m_t>m_Z$), one neglects the fermion masses. Then left and right decays decouple and the partial widths are:
\begin{align}
\boxed{
\Gamma(Z^0\to l_L\bar l_R)=\frac{\alpha_W m_Z}{6c_W^2}\,(C_L^l)^2,
\qquad
\Gamma(Z^0\to l_R\bar l_L)=\frac{\alpha_W m_Z}{6c_W^2}\,(C_R^l)^2,
}
\label{eq:Zleptonic}
\\[1mm]
\boxed{
\Gamma(Z^0\to q_L\bar q_R)=\frac{\alpha_W m_Z N_c}{6c_W^2}\,(C_L^q)^2,
\qquad
\Gamma(Z^0\to q_R\bar q_L)=\frac{\alpha_W m_Z N_c}{6c_W^2}\,(C_R^q)^2.
}
\label{eq:Zhadronic}
\end{align}

\subsection{Total $Z^0$ width and the number of light neutrinos}
Summing over all accessible fermions gives
\begin{equation}
\[
\boxed{
\begin{aligned}
\Gamma(Z^0\to \text{all})
&=
\frac{\alpha_W m_Z}{6c_W^2}\Bigg[
N_\nu\left(\frac{1}{2}\right)^2 \\
&\quad
+ 3\left(\left(-\frac{1}{2}+s_W^2\right)^2+s_W^4\right) \\
&\quad
+ 2N_c\left(\left(\frac{1}{2}-\frac{2s_W^2}{3}\right)^2+\left(-\frac{2s_W^2}{3}\right)^2\right) \\
&\quad
+ 3N_c\left(\left(-\frac{1}{2}+\frac{s_W^2}{3}\right)^2+\left(\frac{s_W^2}{3}\right)^2\right)
\Bigg].
\end{aligned}
}
\]

\label{eq:Ztotal}
\end{equation}
Here $N_\nu$ is the number of neutrinos with mass $m_\nu<m_Z/2$.
Therefore an accurate measurement of $\Gamma(Z^0\to\text{all})$ determines the number of \emph{light} neutrinos, and hence the number of families in the electroweak theory.

Using $m_Z\simeq 91.2~\mathrm{GeV}$, $\alpha_W\simeq 0.0336$ and $s_W^2\simeq 0.23$, the quoted values (in GeV) are:

\begin{center}
\begin{tabular}{cccc}
\hline
$N_\nu$ & EW & EW+QCD & Exp \\
\hline
2 & 2.26 & 2.38 & -- \\
3 & 2.42 & 2.55 & 2.4952(23) \\
4 & 2.59 & 2.71 & -- \\
\hline
\end{tabular}
\end{center}

\noindent
The LEP measurement (00) thus points to $N_\nu=3$ light neutrinos.

\section{Summary of boxed key results}
For quick reference, the most important relations established in this chapter are:
\begin{itemize}
  \item Charge--hypercharge relation:
  \[
  \boxed{Q=T_3+\frac{Y}{2}.}
  \]
  \item Electroweak coupling relation defining the electric charge:
  \[
  \boxed{g\,s_W=g'\,c_W=e.}
  \]
  \item Neutral coupling:
  \[
  \boxed{g_Z=\frac{e}{c_W s_W}.}
  \]
  \item Higgs breaking pattern and gauge boson masses:
  \[
  \boxed{SU_L(2)\times U_Y(1)\to U_{\rm em}(1),\qquad
  m_W=\frac{gv}{2},\quad m_Z=\frac{g_Z v}{2}.}
  \]
  \item Matching to Fermi theory (tree level):
  \[
  \boxed{\frac{g^2}{2m_W^2}=2\sqrt{2}\,G,\quad
  \rho=1,\quad v\simeq 246~\mathrm{GeV}.}
  \]
  \item Fermion masses from Yukawas:
  \[
  \boxed{m_f=\frac{\lambda_f v}{\sqrt{2}},\qquad
  \mathcal L_{h\bar f f}\propto \frac{m_f}{v}\,h.}
  \]
  \item CKM matrix in the charged current:
  \[
  \boxed{\bar u_L\gamma^\mu d_L\ \to\ \bar u_L\gamma^\mu V_{\rm CKM} d_L.}
  \]
  \item Representative electroweak widths:
  \[
  \boxed{\Gamma(W\to l\nu)=\frac{\alpha_W m_W}{12},\qquad
  \Gamma(W\to \text{hadrons})\simeq \frac{\alpha_W m_W N_c}{6}.}
  \]
  \item Total $Z$ width and $N_\nu$:
  \[
  \boxed{\Gamma(Z\to \text{all})\ \text{fixes}\ N_\nu\ \text{(number of light neutrinos)}.}
  \]
\end{itemize}
